\documentclass[11pt]{article}

\usepackage[a4paper,margin=1in]{geometry}
\usepackage{amsmath,amssymb,amsthm,mathtools}
\usepackage{enumitem}

\numberwithin{equation}{section}

\newtheorem{theorem}{Theorem}[section]
\newtheorem{lemma}[theorem]{Lemma}
\newtheorem{proposition}[theorem]{Proposition}
\newtheorem{remark}[theorem]{Remark}

\newcommand{\dd}{\,\mathrm d}
\newcommand{\1}{\mathbf 1}
\newcommand{\E}{\mathbb E}
\newcommand{\R}{\mathbb R}
\newcommand{\cB}{\mathcal B}
\newcommand{\cE}{\mathcal E}
\newcommand{\cF}{\mathcal F}
\newcommand{\Var}{\operatorname{Var}}

\begin{document}
\title{Submission C solution to Problem 5}
\date{}
\maketitle


\begin{abstract}
We prove that the Markov semigroup associated with the Dirichlet stochastic heat
equation with drift $\delta_0(u)$ has at most one invariant probability measure.
In fact, under the hypotheses stated in the question, the invariant measure is unique
and is the Gibbs perturbation of the Brownian bridge measure
\[
        \mu(\dd u)=Z^{-1}\exp\!\left(2\int_0^1 \1_{\{u(\xi)>0\}}\dd \xi\right)\gamma(\dd u),
\]
where $\gamma$ is the law of the standard Brownian bridge on $[0,1]$.
The proof uses three inputs: the reversible Gibbs structure of smooth approximations,
the absolute-continuity theorem of Bounebache--Zambotti for the associated
Dirichlet-form resolvent, and irreducibility of the weighted Ornstein--Uhlenbeck
Dirichlet form.
\end{abstract}

\section{Notation and the candidate invariant measure}

Let
\[
        E:=C_0([0,1]),\qquad H:=L^2(0,1).
\]
We regard $E$ as a Borel subset of $H$. Let $\gamma$ denote the law on $E$ of the
standard Brownian bridge from $0$ to $0$ over $[0,1]$. Equivalently, $\gamma$ is the
centered Gaussian measure on $H$ with covariance $(-\Delta_D)^{-1}$, where
$\Delta_D$ is the Dirichlet Laplacian on $(0,1)$.

Let
\[
        \mathsf H(r):=\1_{\{r>0\}},\qquad
        \Psi(u):=\int_0^1 \mathsf H(u(\xi))\,\dd \xi,\qquad u\in E.
\]
Define
\begin{equation}\label{eq:mu-def}
        \mu(\dd u)
        := Z^{-1}e^{2\Psi(u)}\,\gamma(\dd u),
        \qquad
        Z:=\int_E e^{2\Psi(u)}\,\gamma(\dd u).
\end{equation}
Since $0\leq \Psi\leq 1$, the measures $\mu$ and $\gamma$ are equivalent and their
Radon--Nikodym derivative is bounded above and below by positive constants.

For $n\geq 1$, let
\[
        G_{1/n}(r)=\sqrt{\frac{n}{2\pi}}\,e^{-nr^2/2},\qquad
        B_n(r):=\int_{-\infty}^r G_{1/n}(a)\,\dd a,
\]
and put
\[
        \Psi_n(u):=\int_0^1 B_n(u(\xi))\,\dd \xi,\qquad
        \mu_n(\dd u):=Z_n^{-1}e^{2\Psi_n(u)}\,\gamma(\dd u).
\]
Then $B_n'=G_{1/n}$.

\begin{lemma}\label{lem:mu-n-tv}
The measures $\mu_n$ converge to $\mu$ in total variation.
\end{lemma}

\begin{proof}
For a Brownian bridge $\beta$,
\[
    \E_\gamma\!\left[\int_0^1 \1_{\{\beta(\xi)=0\}}\dd \xi\right]
    =\int_0^1 \gamma(\beta(\xi)=0)\dd \xi=0,
\]
because $\beta(\xi)$ is a non-degenerate centered Gaussian random variable for
$\xi\in(0,1)$, while the endpoints have Lebesgue measure zero. Hence, for
$\gamma$-a.e. $u$, the zero set $\{\xi:u(\xi)=0\}$ has Lebesgue measure zero.
Since $B_n(r)\to \mathsf H(r)$ for every $r\neq0$ and $0\leq B_n\leq1$, dominated
convergence gives $\Psi_n(u)\to\Psi(u)$ for $\gamma$-a.e. $u$. Another application
of dominated convergence gives $Z_n\to Z$ and
\[
        Z_n^{-1}e^{2\Psi_n}\longrightarrow Z^{-1}e^{2\Psi}
        \quad\text{in }L^1(\gamma).
\]
This is exactly total-variation convergence.
\end{proof}

\section{Reversibility of the limiting semigroup}

Let $P_t^n$ be the Markov semigroup for the regularized equation
\[
    \partial_t u^n=\frac12\partial_{xx}u^n+G_{1/n}(u^n)+\dot W
\]
with Dirichlet boundary conditions. The following is standard for gradient-type
stochastic heat equations with smooth drift.

\begin{lemma}\label{lem:approx-reversible}
For each $n$, $\mu_n$ is invariant and reversible for $P_t^n$:
\[
        \int_E f\,P_t^n g\,\dd\mu_n
        =
        \int_E g\,P_t^n f\,\dd\mu_n
\]
for all bounded Borel $f,g:E\to\R$ and all $t\geq0$.
\end{lemma}

\begin{proof}
For cylinder functions $\varphi,\psi$ on $H$, the generator of the regularized
equation is
\[
        L_n\varphi=L_0\varphi+\langle D\Psi_n,D\varphi\rangle_H,
\]
where $L_0$ is the Ornstein--Uhlenbeck generator associated with the linear
Dirichlet stochastic heat equation. Since $\mu_n=Z_n^{-1}e^{2\Psi_n}\gamma$, the
Gaussian integration-by-parts formula gives
\[
        \int \psi\,L_n\varphi\,\dd\mu_n
        =
        -\frac12\int \langle D\varphi,D\psi\rangle_H\,\dd\mu_n
        =
        \int \varphi\,L_n\psi\,\dd\mu_n .
\]
Thus the closure of $L_n$ is self-adjoint in $L^2(\mu_n)$, and the associated
Markov semigroup is symmetric. Taking $\psi\equiv1$ gives invariance. This is the
usual Gibbs-reversibility computation for stochastic quantization; see, for example,
Bounebache--Zambotti \cite[Section 4]{BZ14} for the same calculation in the present
Dirichlet setting.
\end{proof}

We now pass to the singular limit.

\begin{lemma}\label{lem:limit-reversible}
The measure $\mu$ is invariant and reversible for $P_t$. Namely, for all bounded
Borel $f,g:E\to\R$ and all $t\geq0$,
\begin{equation}\label{eq:detailed-balance}
        \int_E f\,P_tg\,\dd\mu
        =
        \int_E g\,P_tf\,\dd\mu .
\end{equation}
\end{lemma}

\begin{proof}
First take $f,g\in C_b(E)$. By the assumed weak convergence of the regularized
solutions,
\[
        P_t^n h(u)\longrightarrow P_t h(u),\qquad u\in E,
\]
for every $h\in C_b(E)$ and every fixed $t>0$. Therefore, using Lemma
\ref{lem:mu-n-tv},
\[
\begin{aligned}
\left|\int fP_t^ng\,\dd\mu_n-\int fP_tg\,\dd\mu\right|
&\leq
\|f\|_\infty\|g\|_\infty\,\|\mu_n-\mu\|_{\mathrm{TV}}        \\
&\quad
+\int |f(u)|\,|P_t^ng(u)-P_tg(u)|\,\mu(\dd u)
\longrightarrow0.
\end{aligned}
\]
Similarly,
\[
        \int gP_t^nf\,\dd\mu_n\longrightarrow \int gP_tf\,\dd\mu .
\]
The reversibility identity for $P_t^n$ therefore passes to the limit and yields
\eqref{eq:detailed-balance} for bounded continuous $f,g$. Since $E$ is Polish and
$\mu$ is a probability measure, a monotone-class argument extends
\eqref{eq:detailed-balance} to all bounded Borel $f,g$. Taking $g\equiv1$ gives
invariance of $\mu$.
\end{proof}

\section{The Dirichlet form and its irreducibility}

Let $\mathcal{FC}_b^1$ denote the smooth bounded cylinder functions on $H$. Define
on $\mathcal{FC}_b^1$ the bilinear form
\begin{equation}\label{eq:dirichlet-form}
        \cE(F,G)
        :=
        \frac12\int_E \langle DF(u),DG(u)\rangle_H\,\mu(\dd u).
\end{equation}
Because $\dd\mu/\dd\gamma$ is bounded above and below by positive constants,
\eqref{eq:dirichlet-form} is closable in $L^2(\mu)$. We denote its closure by
$(\cE,\cF)$.

\begin{lemma}\label{lem:dirichlet-form}
The symmetric semigroup $(P_t)_{t\geq0}$ is the Markov semigroup associated with
$(\cE,\cF)$. Moreover, $(\cE,\cF)$ satisfies a Poincare inequality: there exists
$C<\infty$ such that for every $F\in\cF$,
\begin{equation}\label{eq:poincare}
        \Var_\mu(F)\leq C\,\cE(F,F).
\end{equation}
Consequently, if $F\in L^2(\mu)$ and $P_tF=F$ in $L^2(\mu)$ for all $t\geq0$, then
$F$ is $\mu$-a.s. constant.
\end{lemma}

\begin{proof}
For $f=-2\mathsf H$, Bounebache--Zambotti \cite[Propositions 2.4 and 3.1]{BZ14}
construct the quasi-regular Dirichlet form
\[
        \frac12\int \langle DF,DG\rangle\,\dd\nu_f,
        \qquad
        \nu_f(\dd u)=\frac1{Z_f}\exp\!\left(-\int_0^1 f(u(\xi))\,\dd\xi\right)\gamma(\dd u).
\]
Since $f=-2\mathsf H$, we have $\nu_f=\mu$. For smooth $f$, their local-time
formulation is equivalent, by the occupation-time formula, to
\[
        \partial_t u=\frac12\partial_{xx}u-\frac12 f'(u)+\dot W.
\]
For $f=-2\mathsf H$, one has $-\frac12 f'=\delta_0$ in the sense of distributions.
Thus the Bounebache--Zambotti process solves the same singular SPDE. By the
assumed uniqueness in law of the ABLM solution, the associated semigroup agrees
with $P_t$. Hence $P_t$ is the semigroup of $(\cE,\cF)$.

It remains to prove \eqref{eq:poincare}. Let $\cE_0$ be the Ornstein--Uhlenbeck
Dirichlet form on $L^2(\gamma)$:
\[
        \cE_0(F,F)=\frac12\int \|DF\|_H^2\,\dd\gamma .
\]
The Brownian-bridge Ornstein--Uhlenbeck form has a spectral gap; equivalently,
there exists $C_0<\infty$ such that
\[
        \Var_\gamma(F)\leq C_0\,\cE_0(F,F),
        \qquad F\in\operatorname{Dom}(\cE_0).
\]
This follows from the Hermite expansion of the Dirichlet Ornstein--Uhlenbeck
generator; see Bounebache--Zambotti \cite[Proposition 2.1]{BZ14}.

Since $\dd\mu=m\,\dd\gamma$ with $0<a\leq m\leq b<\infty$, the form domains of
$\cE$ and $\cE_0$ coincide and
\[
        a\,\cE_0(F,F)\leq \cE(F,F)\leq b\,\cE_0(F,F).
\]
Therefore,
\[
\begin{aligned}
        \Var_\mu(F)
        &=\inf_{c\in\R}\int (F-c)^2\,\dd\mu                                      \\
        &\leq \int \left(F-\int F\,\dd\gamma\right)^2\,\dd\mu                       \\
        &\leq b\,\Var_\gamma(F)
        \leq bC_0\,\cE_0(F,F)
        \leq \frac{bC_0}{a}\,\cE(F,F).
\end{aligned}
\]
This proves \eqref{eq:poincare}. If $P_tF=F$ for all $t$, then
\[
        \cE(F,F)
        =
        \lim_{t\downarrow0}\frac1t\langle F-P_tF,F\rangle_{L^2(\mu)}
        =0,
\]
and \eqref{eq:poincare} implies $\Var_\mu(F)=0$.
\end{proof}

We also need the absolute-continuity property of the resolvent. For $\alpha>0$,
write the normalized resolvent as
\[
        \mathcal R_\alpha(x,A)
        :=
        \alpha\int_0^\infty e^{-\alpha t}P_t(x,A)\,\dd t,
        \qquad x\in E,\ A\in\cB(E).
\]

\begin{lemma}[Resolvent absolute continuity]\label{lem:resolvent-ac}
For every $\alpha>0$ and every $x\in E$,
\[
        \mathcal R_\alpha(x,\cdot)\ll \mu.
\]
\end{lemma}

\begin{proof}
Bounebache--Zambotti prove the absolute-continuity condition for the resolvent of
the Dirichlet-form process associated with $\nu_f$ for every bounded function
$f$ of bounded variation; see \cite[Proposition 2.5]{BZ14}. Applying their result
to $f=-2\mathsf H$ gives absolute continuity with respect to $\mu$. Their
Proposition 3.1 identifies the associated process as a weak solution of the
corresponding local-time SPDE. As explained in the proof of Lemma
\ref{lem:dirichlet-form}, for $f=-2\mathsf H$ this is precisely the SPDE with drift
$\delta_0(u)$. The assumed uniqueness in law of the ABLM mild solution therefore
identifies this resolvent with the resolvent of $P_t$. Hence
$\mathcal R_\alpha(x,\cdot)\ll\mu$.
\end{proof}

\section{Uniqueness of the invariant measure}

\begin{theorem}\label{thm:unique-invariant}
The semigroup $(P_t)_{t\geq0}$ has exactly one invariant probability measure on
$E=C_0([0,1])$. It is the measure $\mu$ defined in \eqref{eq:mu-def}. In particular,
$P_t$ has at most one invariant probability measure.
\end{theorem}

\begin{proof}
By Lemma \ref{lem:limit-reversible}, $\mu$ is invariant. Let $\lambda$ be any
invariant probability measure for $P_t$. We first show that $\lambda\ll\mu$.

Let $A\in\cB(E)$ satisfy $\mu(A)=0$. By Lemma \ref{lem:resolvent-ac},
$\mathcal R_\alpha(x,A)=0$ for every $x\in E$. Since $\lambda$ is invariant,
\[
\begin{aligned}
        \lambda(A)
        &=
        \alpha\int_0^\infty e^{-\alpha t}\lambda P_t(A)\,\dd t        \\
        &=
        \int_E \mathcal R_\alpha(x,A)\,\lambda(\dd x)
        =0.
\end{aligned}
\]
Thus $\lambda\ll\mu$. Write $h:=\dd\lambda/\dd\mu$. Then $h\geq0$ and
$\int h\,\dd\mu=1$.

We claim that $P_t h=h$ in $L^1(\mu)$ for every $t\geq0$. Indeed, for every bounded
Borel $\varphi$,
\[
\begin{aligned}
        \int_E \varphi\,P_t h\,\dd\mu
        &=
        \int_E P_t\varphi\, h\,\dd\mu                         \\
        &=
        \int_E P_t\varphi\,\dd\lambda                         \\
        &=
        \int_E \varphi\,\dd\lambda                            \\
        &=
        \int_E \varphi\,h\,\dd\mu .
\end{aligned}
\]
Here the first equality uses the symmetry of $P_t$ in $L^2(\mu)$, extended by
duality to $L^1(\mu)$--$L^\infty(\mu)$, and the third equality uses invariance of
$\lambda$.

Fix $N>0$ and set $h_N:=h\wedge N$. Since $r\mapsto r\wedge N$ is concave and
nondecreasing, Jensen's inequality for the Markov operator $P_t$ gives
\[
        P_t h_N \leq (P_t h)\wedge N=h_N.
\]
Integrating with respect to the invariant measure $\mu$,
\[
        \int P_t h_N\,\dd\mu=\int h_N\,\dd\mu.
\]
Therefore $P_t h_N=h_N$ $\mu$-a.s. for every $t\geq0$. Since $h_N\in L^2(\mu)$,
Lemma \ref{lem:dirichlet-form} implies that $h_N$ is $\mu$-a.s. constant.

Taking $N=2$, we obtain $h\wedge2=c$ for some constant $c$. Since
\[
        c=\int (h\wedge2)\,\dd\mu\leq \int h\,\dd\mu=1<2,
\]
the identity $h\wedge2=c$ forces $h=c$ $\mu$-a.s. Finally,
$\int h\,\dd\mu=1$ gives $c=1$. Hence $\lambda=\mu$.

Thus $\mu$ is the unique invariant probability measure.
\end{proof}

\begin{remark}
The key point excluding additional singular invariant measures is the pointwise
absolute continuity of the normalized resolvent in Lemma \ref{lem:resolvent-ac}.
Symmetry and the Poincare inequality alone imply uniqueness only among invariant
probability measures that are already absolutely continuous with respect to $\mu$.
\end{remark}

\begin{thebibliography}{99}

\bibitem[ABLM24]{ABLM24}
S.~Athreya, O.~Butkovsky, K.~L\^e and L.~Mytnik,
\emph{Well-posedness of stochastic heat equation with distributional drift and skew stochastic heat equation},
Comm. Pure Appl. Math. \textbf{77} (2024), no.~5, 2708--2777.

\bibitem[BZ14]{BZ14}
S.~K.~Bounebache and L.~Zambotti,
\emph{A skew stochastic heat equation},
J. Theoret. Probab. \textbf{27} (2014), no.~1, 168--201.

\bibitem[FOT11]{FOT11}
M.~Fukushima, Y.~Oshima and M.~Takeda,
\emph{Dirichlet Forms and Symmetric Markov Processes},
2nd revised and extended edition, de Gruyter, Berlin, 2011.

\end{thebibliography}

\end{document}